% =========================================================================
%  GREENHERB - Relatório de Projeto
%  Template LaTeX para Overleaf
%  ---------------------------------------------------------------
%  Como usar:
%    1. Substituir os campos em config/identidade.tex (autores, turma, etc.)
%    2. Editar os capítulos em capitulos/
%    3. Adicionar imagens em img/
%    4. Compilar com pdfLaTeX (Overleaf: Menu > Compiler > pdfLaTeX)
% =========================================================================

\documentclass[11pt, a4paper]{report}

% --- Codificação e idioma --------------------------------------
\usepackage[utf8]{inputenc}
\usepackage[T1]{fontenc}
\usepackage[portuguese]{babel}
\usepackage{lmodern}

% --- Geometria da página ---------------------------------------
\usepackage[a4paper, top=2.5cm, bottom=2.5cm, left=2.8cm, right=2.8cm, headheight=14pt]{geometry}

% --- Tipografia ------------------------------------------------
\usepackage{microtype}
\usepackage{setspace}
\onehalfspacing
\usepackage{parskip}     % parágrafos com espaço em vez de indentação

% --- Cores (paleta Forest & Moss, alinhada com os slides) ------
\usepackage{xcolor}
\definecolor{forest}{HTML}{2C5F2D}
\definecolor{moss}{HTML}{97BC62}
\definecolor{mossLight}{HTML}{C8DBA0}
\definecolor{charcoal}{HTML}{1F2D1F}
\definecolor{textMuted}{HTML}{5C6B5C}
\definecolor{amber}{HTML}{D9822B}
\definecolor{boxbg}{HTML}{F5F8F0}
\definecolor{border}{HTML}{D8DDD0}

% --- Gráficos e tabelas ----------------------------------------
\usepackage{graphicx}
\usepackage{float}
\usepackage{tabularx}
\usepackage{booktabs}
\usepackage{array}
\usepackage{longtable}
\usepackage{multirow}
\usepackage{amssymb}     % \checkmark e símbolos matemáticos

% --- Listas e código -------------------------------------------
\usepackage{enumitem}
\setlist[itemize]{topsep=2pt, itemsep=1pt, parsep=0pt}
\setlist[enumerate]{topsep=2pt, itemsep=1pt, parsep=0pt}

\usepackage{listings}
\lstset{
  basicstyle=\ttfamily\small,
  backgroundcolor=\color{boxbg},
  frame=single,
  rulecolor=\color{border},
  framesep=4pt,
  breaklines=true,
  showstringspaces=false,
  captionpos=b,
  keywordstyle=\color{forest}\bfseries,
  commentstyle=\color{textMuted}\itshape,
  stringstyle=\color{amber}
}

% --- Caixas destacadas (notas, dicas, decisões) ----------------
\usepackage[most]{tcolorbox}

\newtcolorbox{decisao}[1][]{
  colback=boxbg, colframe=forest, boxrule=0.6pt,
  arc=2pt, left=10pt, right=10pt, top=8pt, bottom=8pt,
  fonttitle=\bfseries\color{white},
  coltitle=white,
  colbacktitle=forest,
  title=Decisão de arquitetura,
  #1
}

\newtcolorbox{nota}[1][]{
  colback=mossLight!30, colframe=moss, boxrule=0.6pt,
  arc=2pt, left=10pt, right=10pt, top=6pt, bottom=6pt,
  #1
}

\newtcolorbox{aviso}[1][]{
  colback=amber!10, colframe=amber, boxrule=0.6pt,
  arc=2pt, left=10pt, right=10pt, top=6pt, bottom=6pt,
  #1
}

% --- Diagramas (TikZ) ------------------------------------------
\usepackage{tikz}
\usetikzlibrary{shapes.geometric, arrows.meta, positioning, fit, backgrounds}

% --- Hyperlinks ------------------------------------------------
\usepackage[hidelinks, colorlinks=true, linkcolor=forest, urlcolor=forest, citecolor=forest]{hyperref}
\usepackage{url}

% --- Cabeçalhos e rodapés --------------------------------------
\usepackage{fancyhdr}
\pagestyle{fancy}
\fancyhf{}
\fancyhead[L]{\small\color{textMuted}\nouppercase{\leftmark}}
\fancyhead[R]{\small\color{textMuted}GREENHERB}
\fancyfoot[C]{\thepage}
\renewcommand{\headrulewidth}{0.4pt}
\renewcommand{\headrule}{\hbox to\headwidth{\color{border}\leaders\hrule height \headrulewidth\hfill}}

% --- Estilo dos títulos de capítulo ----------------------------
\usepackage{titlesec}

\titleformat{\chapter}[display]
  {\normalfont\sffamily\Huge\bfseries\color{forest}}
  {\fontsize{60}{60}\selectfont\color{moss}\thechapter}
  {0pt}
  {\titlerule[0.6pt]\vspace{6pt}\Huge}

\titleformat{\section}
  {\normalfont\sffamily\Large\bfseries\color{forest}}
  {\thesection}{1em}{}

\titleformat{\subsection}
  {\normalfont\sffamily\large\bfseries\color{charcoal}}
  {\thesubsection}{1em}{}

\titleformat{\subsubsection}
  {\normalfont\sffamily\normalsize\bfseries\color{charcoal}}
  {\thesubsubsection}{1em}{}

\titlespacing*{\chapter}{0pt}{-30pt}{30pt}

% --- Bibliografia ----------------------------------------------
\usepackage[backend=biber, style=numeric, sorting=none]{biblatex}
\addbibresource{bibliografia.bib}

% --- Identidade (autores, instituição, etc.) -------------------
% =========================================================================
%  Identidade do trabalho
%  ---------------------------------------------------------------
%  Os alunos devem editar APENAS este ficheiro para configurar o
%  cabeçalho, autores e contexto institucional do relatório.
% =========================================================================

% --- Trabalho --------------------------------------------------
\newcommand{\projetonome}{GREENHERB}
\newcommand{\projetosubtitulo}{Plataforma Web de Gestão Inteligente de uma Estufa}
\newcommand{\relatoriotipo}{Relatório de Projeto}

% --- Instituição -----------------------------------------------
\newcommand{\instituicao}{Universidade da Beira Interior}
\newcommand{\curso}{Licenciatura em ...}
\newcommand{\unidadecurricular}{Unidade Curricular de ...}
\newcommand{\anoletivo}{2025/2026}
\newcommand{\semestreletivo}{Semestre 2}

% --- Equipa ----------------------------------------------------
% Adicionar / remover linhas conforme o número de elementos
\newcommand{\equipa}{%
  \begin{tabular}{ll}
    Nome do Aluno 1 & Nº 99999 \\
    Nome do Aluno 2 & Nº 99999 \\
    Nome do Aluno 3 & Nº 99999 \\
  \end{tabular}%
}

% --- Docente ---------------------------------------------------
\newcommand{\docente}{Nome do/a Docente}

% --- Data ------------------------------------------------------
\newcommand{\datarelatorio}{\today}

% --- Repositório (opcional) ------------------------------------
\newcommand{\repositorio}{https://github.com/equipa/greenherb}


% =========================================================================
\begin{document}

% --- Capa ------------------------------------------------------
% =========================================================================
%  Capa
% =========================================================================

\begin{titlepage}
\thispagestyle{empty}

% Banda decorativa superior
\begin{tikzpicture}[remember picture, overlay]
  \fill[forest] (current page.north west) rectangle ([yshift=-3cm]current page.north east);

  % Pequeno acento moss
  \fill[moss] ([yshift=-3cm]current page.north west) rectangle ([yshift=-3.15cm]current page.north east);

  % Texto da instituição centrado dentro da banda
  \node[anchor=center, text=white, align=center]
    at ([yshift=-1.5cm]current page.north) {%
      {\sffamily\Large\bfseries \instituicao}
    };
\end{tikzpicture}

\vspace*{4cm}

% Título principal
\begin{center}
  {\sffamily\color{moss}\large \relatoriotipo}\\[1.5cm]
  {\sffamily\Huge\bfseries\color{forest} \projetonome}\\[0.6cm]
  {\sffamily\large\itshape\color{charcoal} \projetosubtitulo}
\end{center}

\vspace{1.5cm}

\begin{center}
  \textcolor{moss}{\rule{4cm}{1pt}}
\end{center}

\vspace{1.5cm}

% Equipa
\begin{center}
  {\sffamily\color{textMuted}\small REALIZADO POR}\\[0.4cm]
  {\sffamily\normalsize \equipa}
\end{center}

\vfill

% Rodapé com info do curso
\begin{center}
  \begin{tabular}{c}
    {\sffamily\small\color{charcoal} \curso} \\[0.2cm]
    {\sffamily\small\color{textMuted} \unidadecurricular} \\[0.2cm]
    {\sffamily\small\color{textMuted} Docente: \docente} \\[0.4cm]
    {\sffamily\small\color{textMuted} \anoletivo \quad $\cdot$ \quad \semestreletivo} \\[0.2cm]
    {\sffamily\small\color{textMuted} \datarelatorio}
  \end{tabular}
\end{center}

\vspace{0.5cm}

\end{titlepage}


% --- Páginas iniciais ------------------------------------------
\pagenumbering{roman}
\setcounter{page}{1}

% =========================================================================
%  Resumo
% =========================================================================

\chapter*{Resumo}
\addcontentsline{toc}{chapter}{Resumo}

\noindent
% --- INSTRUÇÕES (apagar antes da entrega) ---------------------
% O resumo deve ter entre 150 e 250 palavras e responder a:
%   1. Qual o problema que o projeto resolve?
%   2. Qual a abordagem técnica escolhida?
%   3. Quais os principais resultados ou contribuições?
% Evitar entrar em detalhes técnicos extensos — esses ficam para
% os capítulos seguintes.
% ---------------------------------------------------------------

[Substituir pelo texto do resumo. Apresentar em poucas linhas o
contexto do projeto GREENHERB, os principais requisitos abordados, a
arquitetura geral adotada (browser, API, base de dados) e os resultados
relevantes. O leitor que apenas leia este resumo deve ficar com uma
ideia clara do que foi feito e do que se demonstrou.]

\vspace{0.5cm}
\noindent\textbf{Palavras-chave:} aplicação web, gestão de estufa, Node.js,
MongoDB, armazenamento no browser, JWT, OpenAPI.

\clearpage


\tableofcontents
\clearpage

\listoffigures
\listoftables
\clearpage

% --- Corpo principal -------------------------------------------
\pagenumbering{arabic}
\setcounter{page}{1}

% =========================================================================
%  Capítulo 1 — Introdução
% =========================================================================

\chapter{Introdução}
\label{cap:introducao}

% --- INSTRUÇÕES (apagar antes da entrega) ---------------------
% A introdução deve permitir a um leitor externo perceber:
%   - O problema que o sistema vem resolver.
%   - O contexto em que o sistema vai ser utilizado.
%   - Quais os objetivos do projeto.
%   - Como está organizado o restante do relatório.
% ---------------------------------------------------------------

\section{Enquadramento}
\label{sec:enquadramento}

% Apresentem aqui o domínio do problema. O que é a GREENHERB?
% Que tipo de operação se passa numa estufa de ervas aromáticas?
% Que dificuldades existem hoje sem informatização?
% (1 a 2 parágrafos)

A empresa GREENHERB pretende informatizar a gestão de uma estufa
dedicada à produção de ervas aromáticas, designadamente
manjericão, hortelã, alecrim e tomilho. O sistema atual,
maioritariamente manual, cria dificuldades no planeamento, na
monitorização das condições ambientais e na rastreabilidade das
intervenções realizadas.

[Completar: descrever o impacto destas dificuldades na operação
quotidiana da estufa e a oportunidade que a digitalização representa.]


\section{Objetivos}
\label{sec:objetivos}

% Listar os objetivos do projeto. Recomenda-se uma lista
% numerada, com objetivos verificáveis (não vagos).

O presente trabalho tem como objetivos:

\begin{enumerate}
  \item Desenvolver uma aplicação web \emph{full-stack} que suporte o ciclo de vida completo de lotes de cultivo;
  \item Implementar autenticação, autorização baseada em perfis e rastreabilidade das ações realizadas;
  \item Especificar e documentar a API segundo a norma OpenAPI 3.x;
  \item Definir uma arquitetura de armazenamento no browser que permita o funcionamento da aplicação em condições de conectividade instável;
  \item {}[Adicionar / ajustar conforme decisões da equipa.]
\end{enumerate}


\section{Contexto e restrições}
\label{sec:contexto}

% Esta secção é onde se reconhecem as restrições do projeto.
% Útil para enquadrar decisões posteriores.

A solução é desenvolvida com restrições técnicas previamente
definidas: o \emph{frontend} em \texttt{HTML}, \texttt{CSS} e
\texttt{JavaScript}, o \emph{backend} em \texttt{Node.js} com
recurso a \texttt{Mongoose} para acesso a uma base de dados
\texttt{MongoDB}, autenticação baseada em \texttt{JWT} e uso obrigatório
dos mecanismos de armazenamento do navegador (\texttt{localStorage}
ou \texttt{sessionStorage}, \texttt{IndexedDB} e \texttt{Cache API}).

Acresce o contexto operacional: a aplicação é utilizada em ambiente
de estufa onde a conectividade pode ser instável, o que condiciona
fortemente as decisões de arquitetura do lado do cliente.


\section{Estrutura do documento}
\label{sec:estrutura}

% Pequena ``mapa'' do que o leitor vai encontrar nos próximos
% capítulos. Manter curto.

O restante documento organiza-se da seguinte forma. O Capítulo~\ref{cap:requisitos}
identifica e especifica os requisitos do sistema, distinguindo entre
funcionais e não funcionais e apresentando a sua priorização. O
Capítulo~\ref{cap:arquitetura} descreve a arquitetura adotada,
detalhando as responsabilidades de cada camada (\emph{browser}, API e
base de dados) e justificando as decisões tomadas. O
Capítulo~\ref{cap:conclusao} apresenta as conclusões, balanço crítico
e trabalho futuro.

% =========================================================================
%  Capítulo 2 — Requisitos
% =========================================================================

\chapter{Requisitos}
\label{cap:requisitos}

% --- INSTRUÇÕES (apagar antes da entrega) ---------------------
% Este capítulo deve identificar e especificar os requisitos do
% sistema. Recomenda-se a divisão em:
%   - Atores e perfis de utilizador
%   - Regras de negócio relevantes
%   - Requisitos funcionais (com identificador único e descrição)
%   - Requisitos não funcionais
%   - Priorização e roadmap por sprint
% Cada requisito deve ser RASTREÁVEL: ter um ID, uma descrição
% verificável e, idealmente, um critério de aceitação.
% ---------------------------------------------------------------


\section{Atores do sistema}
\label{sec:atores}

% Identificar os perfis de utilizador e o que cada um pode fazer.
% Aproveitar a tabela abaixo como ponto de partida.

O sistema considera três perfis de utilizador, cada um com um
âmbito distinto de operações.

\begin{table}[H]
\centering
\caption{Perfis de utilizador e respetivas responsabilidades.}
\label{tab:atores}
\renewcommand{\arraystretch}{1.3}
\begin{tabularx}{\linewidth}{l X}
\toprule
\textbf{Perfil} & \textbf{Responsabilidades principais} \\
\midrule
Técnico        & Registo de medições e tarefas operacionais; consulta de planos e lotes atribuídos. \\
Responsável    & Definição de planos de cultivo; autorização de intervenções pontuais; resolução de alertas. \\
Administrador  & Gestão de utilizadores e configuração geral do sistema; acesso integral aos dados. \\
\bottomrule
\end{tabularx}
\end{table}


\section{Regras de negócio}
\label{sec:regras-negocio}

% Listar as regras de negócio que condicionam os requisitos.
% É boa prática numerar (RN-01, RN-02, ...) para depois poder
% referenciar nos requisitos e nos critérios de aceitação.

As regras de negócio que condicionam o comportamento do sistema
estão sintetizadas a seguir.

\begin{description}[style=nextline, leftmargin=2cm, labelindent=0pt]
  \item[\textbf{RN-01}] Cada plano de cultivo está associado a uma erva aromática e classifica-se como \emph{regular}, \emph{emergência} ou \emph{pontual}.
  \item[\textbf{RN-02}] Um plano \emph{regular} inclui obrigatoriamente intervalos de temperatura, humidade e luminosidade, plano de rega e fertilização, e duração prevista do ciclo.
  \item[\textbf{RN-03}] Um plano \emph{pontual} exige autorização explícita do Responsável.
  \item[\textbf{RN-04}] Os alertas classificam-se como \emph{Informativo}, \emph{Aviso} ou \emph{Crítico} e podem ser \emph{Resolvidos} ou \emph{Ignorados}; a opção \emph{Ignorado} exige justificação.
  \item[\textbf{RN-05}] O sistema opera em modo \emph{Manual} ou \emph{Automático}; em modo automático pode sugerir ou executar ações com base em regras.
  \item[\textbf{RN-06}] [Adicionar as restantes regras relevantes derivadas do enunciado.]
\end{description}


\section{Priorizacao e roadmap por sprint}
\label{sec:roadmap-sprints}

A equipa adota uma abordagem incremental orientada ao contrato de API,
de forma a reduzir retrabalho entre implementação, documentação e
avaliação. Na Sprint 1 foi concluída a especificação OpenAPI dos
principais recursos de domínio, sem implementação do fluxo de
autenticação.

\begin{table}[H]
\centering
\caption{Roadmap inicial de sprints.}
\label{tab:roadmap}
\renewcommand{\arraystretch}{1.3}
\begin{tabularx}{\linewidth}{l l X}
\toprule
\textbf{Sprint} & \textbf{Foco} & \textbf{Entregáveis principais} \\
\midrule
Sprint 1 & Especificação OpenAPI & \texttt{openapi.yaml}, matriz RF$\rightarrow$endpoints, validação Swagger e alinhamento do relatório. \\
Sprint 2 & Base backend e autenticação & Implementação Node.js/Mongoose, JWT, controlo de acesso por perfil e validações de entrada. \\
Sprint 3 & Regras de negócio e histórico & Medições, alertas, modos manual/automático, histórico e auditoria. \\
Sprint 4 & Offline, exportação e consolidação & Sincronização degradada, exportação CSV/Excel, testes finais e fecho dos critérios de aceitação. \\
\bottomrule
\end{tabularx}
\end{table}


\section{Requisitos funcionais}
\label{sec:rf}

% Cada requisito funcional deve ter:
%   - Um ID único (RF-XX).
%   - Uma descrição clara, na forma "O sistema deve...".
%   - Os atores envolvidos.
%   - Uma prioridade (Alta / Média / Baixa) ou MoSCoW
%     (Must / Should / Could / Won't).
%   - Idealmente, um critério de aceitação verificável.

Os requisitos funcionais identificados estão listados na
Tabela~\ref{tab:rf}. A prioridade segue a notação MoSCoW: \textbf{M}
(\emph{must}, obrigatório), \textbf{S} (\emph{should}, importante),
\textbf{C} (\emph{could}, desejável), \textbf{W} (\emph{won't},
fora do âmbito desta versão).

\begin{table}[H]
\centering
\caption{Lista priorizada de requisitos funcionais.}
\label{tab:rf}
\small
\renewcommand{\arraystretch}{1.3}
\begin{tabularx}{\linewidth}{l X c c}
\toprule
\textbf{ID} & \textbf{Descrição} & \textbf{Atores} & \textbf{Prio.} \\
\midrule
RF-01 & Autenticação de utilizadores e controlo de acesso por perfil. & Todos & M \\
RF-02 & Gestão de utilizadores (criar, editar, desativar) pelo Administrador. & Admin. & M \\
RF-03 & Registo, edição, consulta e importação (CSV/Excel) de ervas aromáticas. & Resp. & M \\
RF-04 & Criação e gestão de planos de cultivo (\emph{regular}, \emph{emergência}, \emph{pontual}). & Resp. & M \\
RF-05 & Associação de planos a lotes e consulta do estado atual de cada lote. & Resp., Téc. & M \\
RF-06 & Registo de tarefas operacionais e marcação como executadas. & Téc. & M \\
RF-07 & Registo de medições ambientais (temperatura, humidade, luminosidade). & Téc. & M \\
RF-08 & Geração automática de alertas quando os valores violam o plano. & Sistema & M \\
RF-09 & Resolução ou \emph{ignore} de alertas, com justificação obrigatória quando aplicável. & Resp. & M \\
RF-10 & Suporte aos modos de operação \emph{Manual} e \emph{Automático}. & Resp. & S \\
RF-11 & Histórico de medições, intervenções e alertas. & Todos & S \\
RF-12 & Comparação entre valores reais e valores esperados pelo plano. & Resp. & S \\
RF-13 & Exportação de relatórios em formato CSV ou Excel. & Resp., Admin. & S \\
RF-14 & Divisão de lotes, registo de perdas e cálculo de produtividade. & Resp. & S \\
RF-15 & Registo de logs de auditoria sobre as operações relevantes. & Sistema & M \\
\bottomrule
\end{tabularx}
\end{table}


\subsection{Detalhe e critérios de aceitação (exemplos)}
\label{sec:rf-detalhe}

% Para os requisitos mais críticos, vale a pena escrever um
% pequeno detalhe com critérios de aceitação. Aqui ficam dois
% exemplos para a equipa replicar nos seus próprios requisitos
% críticos.

A seguir detalham-se dois requisitos a título de exemplo. Recomenda-se
que esta especificação seja replicada para todos os requisitos
classificados como \textbf{M}.

\paragraph{RF-08 — Geração automática de alertas.}
Quando uma medição ambiental é registada e algum dos seus valores
está fora do intervalo definido no plano de cultivo associado ao
lote, o sistema deve criar um alerta com a classificação adequada
(\emph{Informativo}, \emph{Aviso} ou \emph{Crítico}).

\textbf{Critérios de aceitação:}
\begin{itemize}
  \item Dada uma medição cujo valor de temperatura excede em mais de 5\,\%${}^{\circ}$C o limite superior do plano, é gerado um alerta \emph{Crítico} no prazo de 5 segundos.
  \item Dado um sensor que não envia dados há mais de 10 minutos, é gerado um alerta de \emph{Aviso} associado ao lote.
  \item Os alertas gerados ficam visíveis no dashboard do Responsável e no histórico do lote.
\end{itemize}

\paragraph{RF-09 — Resolução ou \emph{ignore} de alertas.}
O Responsável deve poder marcar um alerta como \emph{Resolvido} ou
\emph{Ignorado}. No caso de \emph{ignore}, é obrigatório fornecer
uma justificação textual.

\textbf{Critérios de aceitação:}
\begin{itemize}
  \item A interface impede submeter o \emph{ignore} sem justificação.
  \item A decisão é registada no log de auditoria com identificação do utilizador, momento e justificação.
  \item Um alerta resolvido ou ignorado deixa de aparecer na lista ativa, mas mantém-se no histórico.
\end{itemize}


\section{Requisitos não funcionais}
\label{sec:rnf}

% Estes requisitos não descrevem O QUE o sistema faz, mas COMO
% o faz: desempenho, segurança, usabilidade, manutenibilidade...
% São muitas vezes esquecidos pelos alunos — não os descurem.

\begin{table}[H]
\centering
\caption{Requisitos não funcionais.}
\label{tab:rnf}
\renewcommand{\arraystretch}{1.3}
\begin{tabularx}{\linewidth}{l l X}
\toprule
\textbf{ID} & \textbf{Categoria} & \textbf{Descrição} \\
\midrule
RNF-01 & Segurança       & Palavras-passe armazenadas com \emph{hashing} (e.g., bcrypt); nunca em texto claro. \\
RNF-02 & Segurança       & Comunicações cliente-servidor protegidas por HTTPS em produção. \\
RNF-03 & Segurança       & Autenticação via JWT, com expiração explícita e proteção das rotas. \\
RNF-04 & Disponibilidade & A aplicação deve permitir registo de medições e tarefas em modo degradado (sem rede). \\
RNF-05 & Usabilidade     & Interface funcional em \emph{desktop} e \emph{mobile} (\emph{responsive}). \\
RNF-06 & Manutenibilidade & API documentada em OpenAPI 3.x e alinhada com a implementação. \\
RNF-07 & Auditabilidade  & Todas as operações sensíveis registam autor, momento e ação. \\
RNF-08 & [Adicionar]     & [Outros requisitos relevantes: desempenho, internacionalização, etc.] \\
\bottomrule
\end{tabularx}
\end{table}




% =========================================================================
%  Capítulo 3 — Arquitetura
% =========================================================================

\chapter{Arquitetura}
\label{cap:arquitetura}

% --- INSTRUÇÕES (apagar antes da entrega) ---------------------
% Este capítulo é o que distingue um relatório bom de um
% relatório medíocre. Não basta dizer ``usámos Node.js e
% MongoDB'' — é preciso JUSTIFICAR cada decisão.
%
% Para cada decisão de arquitetura, respondam a três perguntas:
%   1. Que dados/lógica vivem aqui?
%   2. Porquê AQUI e não noutra camada?
%   3. Que alternativas foram consideradas e descartadas?
%
% Usem a caixa "decisao" sempre que registarem uma decisão
% importante — facilita a leitura e mostra o pensamento crítico
% da equipa.
% ---------------------------------------------------------------


\section{Visão geral}
\label{sec:arquitetura-geral}

% Visão de alto nível: três camadas, comunicação entre elas,
% protocolos. Acompanhar com o diagrama abaixo.

A solução adopta uma arquitetura em três camadas, separando
claramente as responsabilidades do cliente, do servidor e da
persistência. A Figura~\ref{fig:arquitetura-geral} apresenta a
visão de alto nível.

\begin{figure}[H]
\centering
\begin{tikzpicture}[
    node distance=1.6cm,
    bloco/.style={
      rectangle, rounded corners=3pt,
      draw=forest, line width=0.8pt,
      fill=boxbg,
      minimum width=3.4cm, minimum height=4cm,
      align=center, font=\small
    },
    seta/.style={
      <->, >=Stealth, line width=0.8pt, color=forest
    }
]

% Bloco BROWSER
\node[bloco] (browser) {%
  \begin{minipage}{3cm}
    \centering
    {\sffamily\bfseries\small\color{forest} BROWSER}\\[1pt]
    {\tiny\itshape\color{textMuted} Cliente}\\[6pt]
    {\scriptsize
      \texttt{HTML} / \texttt{CSS} / \texttt{JS}\\
      \texttt{localStorage}\\
      \texttt{sessionStorage}\\
      \texttt{IndexedDB}\\
      \texttt{Cache API}\\
      Service Worker
    }
  \end{minipage}
};

% Bloco API
\node[bloco, right=of browser] (api) {%
  \begin{minipage}{3cm}
    \centering
    {\sffamily\bfseries\small\color{forest} API NODE.JS}\\[1pt]
    {\tiny\itshape\color{textMuted} Servidor}\\[6pt]
    {\scriptsize
      Rotas REST\\
      Middleware JWT\\
      Validação Mongoose\\
      Motor de regras\\
      Logs auditoria\\
      OpenAPI
    }
  \end{minipage}
};

% Bloco BD
\node[bloco, right=of api] (db) {%
  \begin{minipage}{3cm}
    \centering
    {\sffamily\bfseries\small\color{forest} MONGODB}\\[1pt]
    {\tiny\itshape\color{textMuted} Persistência}\\[6pt]
    {\scriptsize
      Utilizador\\
      ErvaAromática\\
      PlanoCultivo\\
      LoteCultivo\\
      Tarefa, Medição\\
      Alerta, Auditoria
    }
  \end{minipage}
};

% Setas
\draw[seta] (browser) -- node[above, font=\tiny\sffamily, color=textMuted] {HTTPS · JSON} (api);
\draw[seta] (api) -- node[above, font=\tiny\sffamily, color=textMuted] {Mongoose} (db);

\end{tikzpicture}
\caption{Visão de alto nível da arquitetura em três camadas.}
\label{fig:arquitetura-geral}
\end{figure}

A comunicação entre o cliente e o servidor faz-se sobre HTTPS,
trocando mensagens em JSON. O servidor acede à base de dados
através do ODM \texttt{Mongoose}, que assegura a validação de
esquemas antes da persistência.


\section{Matriz de responsabilidades}
\label{sec:matriz-responsabilidades}

% Esta é a tabela mais importante do capítulo. Para cada
% funcionalidade, indica claramente o que fica em cada camada.
% Forçam-se assim respostas explícitas à pergunta "porquê aqui?"

A Tabela~\ref{tab:matriz} apresenta a distribuição de
responsabilidades por camada para as funcionalidades-chave do
sistema.

\begin{table}[H]
\centering
\caption{Distribuição de responsabilidades por camada.}
\label{tab:matriz}
\footnotesize
\renewcommand{\arraystretch}{1.4}
\begin{tabularx}{\linewidth}{l X X X}
\toprule
\textbf{Funcionalidade} & \textbf{Browser} & \textbf{API} & \textbf{Base de Dados} \\
\midrule
Autenticação JWT &
Envio de credenciais; armazenamento e renovação do token. &
Validação de credenciais; emissão e verificação do JWT. &
Persistência de utilizadores e \emph{hashes} de palavra-passe. \\

Registo de medição ambiental &
Captura via formulário; \emph{enqueue} em IndexedDB se \emph{offline}. &
Validação de regras; persistência; verificação de limites. &
Persistência da medição e dos alertas eventualmente gerados. \\

Geração de alertas &
Apresentação ao utilizador; UX de resolução. &
Avaliação dos limites do plano; criação dos alertas. &
Persistência do alerta com referência ao lote e ao plano. \\

Importação de CSV/Excel &
Pré-validação leve do ficheiro (formato, colunas). &
Validação completa, normalização e inserção em \emph{bulk}. &
Persistência das ervas aromáticas importadas. \\

Modo Manual / Automático &
UI de comutação e de listagem das ações sugeridas. &
Motor de regras; execução ou sugestão das ações. &
Registo da configuração do modo e do histórico de ações. \\

Sincronização \emph{offline} &
Fila de operações pendentes em IndexedDB; reenvio ao reconectar. &
Tratamento idempotente das operações; resolução de conflitos. &
Persistência idempotente; carimbos temporais. \\
\bottomrule
\end{tabularx}
\end{table}


\section{Camada de cliente (\emph{browser})}
\label{sec:camada-browser}

% Esta secção é o "documento de arquitetura cliente" exigido
% pelo enunciado. Tem de cobrir os 4 mecanismos de armazenamento.

A camada de cliente foi desenhada para suportar três cenários
distintos de operação: ligação estável, ligação intermitente e
ausência total de conectividade. Esta secção descreve os
mecanismos de armazenamento utilizados, a respetiva justificação,
política de invalidação e estratégia de sincronização.

\subsection{Mecanismos de armazenamento}

\begin{table}[H]
\centering
\caption{Mecanismos de armazenamento no \emph{browser} e respetivos propósitos.}
\label{tab:browser-storage}
\footnotesize
\renewcommand{\arraystretch}{1.4}
\begin{tabularx}{\linewidth}{l X X X}
\toprule
\textbf{Mecanismo} & \textbf{Dados guardados} & \textbf{Justificação} & \textbf{Expiração} \\
\midrule
\texttt{localStorage} &
Preferências da UI; identificador do último ecrã visitado; \emph{feature flags}. &
Pequenos pares chave/valor síncronos; persistem entre sessões. &
Manual; limpa no \emph{logout}. \\

\texttt{sessionStorage} &
Estado temporário de formulários multi-passo (\emph{wizards}). &
Volátil por aba; evita ``poluir'' o estado entre sessões. &
Fim da sessão da aba. \\

\texttt{IndexedDB} &
Histórico recente; cache de lotes e medições; fila de operações pendentes (\emph{offline}). &
Estrutura, queryable, assíncrono; permite operação degradada. &
TTL de [N] dias para o cache; a fila é purgada após sincronização bem-sucedida. \\

\texttt{Cache API} &
Recursos estáticos (HTML, CSS, JS, ícones); respostas \texttt{GET} idempotentes selecionadas. &
Gerida pelo Service Worker; reduz latência e suporta arranque \emph{offline}. &
Estratégia \emph{stale-while-revalidate}; invalidação por versão. \\
\bottomrule
\end{tabularx}
\end{table}


\subsection{Decisão sobre o armazenamento do JWT}
\label{sec:decisao-jwt}

% Aqui é OBRIGATÓRIO justificar com profundidade. O enunciado
% pede explicitamente esta decisão e a sua justificação de
% segurança. Substituir o exemplo abaixo pela decisão real da
% equipa.

\begin{decisao}[title=Decisão DA-01: armazenamento do token JWT]
\textbf{Decisão.} O token JWT é guardado em [\texttt{localStorage} / \texttt{sessionStorage} / \emph{cookie httpOnly}].

\textbf{Alternativas consideradas.}
\begin{itemize}
  \item \texttt{localStorage}: persistência longa, mas exposto a ataques XSS.
  \item \texttt{sessionStorage}: limita a janela de exposição mas obriga a re-autenticar a cada sessão.
  \item \emph{Cookie httpOnly + Secure + SameSite}: imune a XSS por JavaScript mas requer mitigação de CSRF.
\end{itemize}

\textbf{Justificação.} [Justificar com base no perfil de risco do
sistema: estufa em ambiente operacional, sessões longas, baixa
exposição a CSRF se a API não usa formulários tradicionais, etc.]

\textbf{Mitigações.} [Lista de medidas adotadas: \emph{Content Security
Policy} restritiva, \emph{escape} de saídas, sanitização de \emph{inputs},
expiração curta e renovação por \emph{refresh token}, etc.]
\end{decisao}


\subsection{Estratégia de sincronização}
\label{sec:sync}

% Como funciona o offline → online? Como se evitam duplicações?
% Como se resolvem conflitos?

A operação em modo degradado segue o padrão \emph{outbox}: as
operações realizadas \emph{offline} são enfileiradas em \texttt{IndexedDB}
e reenviadas ao servidor assim que a conectividade é restabelecida.
Para garantir a idempotência, cada operação carrega um identificador
único gerado no cliente (UUID) que o servidor reconhece e ignora em
caso de submissões repetidas.

\begin{aviso}
Em caso de conflito (e.g., duas equipas registam a mesma colheita
do mesmo lote em momentos diferentes), o sistema adota a estratégia
de [\emph{last-write-wins} / \emph{merge} / \emph{flag para revisão
humana}]. Esta decisão tem implicações operacionais e deve ser
discutida explicitamente.
\end{aviso}


\section{Camada de servidor (API)}
\label{sec:camada-api}

% Estrutura da API: organização das rotas, middleware,
% autenticação, motor de regras, geração de logs, geração da
% documentação OpenAPI.

\subsection{Organização das rotas}

A API expõe os recursos da aplicação seguindo as convenções REST.
A Tabela~\ref{tab:rotas} sintetiza os principais \emph{endpoints}
agrupados por recurso. A especificação completa encontra-se em
\texttt{openapi.yaml} (ver Anexo~\ref{anexo:openapi}).
Na Sprint 1 foi priorizada a especificação dos recursos de domínio,
ficando a implementação da autenticação para sprint posterior.

\begin{table}[H]
\centering
\caption{Resumo dos principais \emph{endpoints} da API.}
\label{tab:rotas}
\footnotesize
\renewcommand{\arraystretch}{1.3}
\begin{tabularx}{\linewidth}{l l X}
\toprule
\textbf{Método} & \textbf{Rota} & \textbf{Descrição} \\
\midrule
GET   & \texttt{/api/v1/users}                  & Listagem de utilizadores (Admin.). \\
POST  & \texttt{/api/v1/herbs/import}           & Importação em \emph{bulk} (CSV/Excel). \\
POST  & \texttt{/api/v1/cultivation-plans}      & Criação de plano de cultivo (\emph{regular}, \emph{emergência}, \emph{pontual}). \\
POST  & \texttt{/api/v1/measurements}           & Registo de medição ambiental com potencial geração de alerta. \\
POST  & \texttt{/api/v1/alerts/:id/ignore}      & Ignorar alerta com justificação obrigatória. \\
GET   & \texttt{/api/v1/reports/export}         & Exportação de relatório em CSV/Excel. \\
\bottomrule
\end{tabularx}
\end{table}


\subsection{Cadeia de \emph{middleware}}

Cada pedido atravessa uma cadeia de \emph{middleware} ordenada da
seguinte forma:

\begin{enumerate}
  \item \textbf{Logger HTTP}: regista método, rota, IP e \emph{user-agent}.
  \item \textbf{Validação JWT}: verifica assinatura e expiração; rejeita pedidos não autenticados em rotas protegidas.
  \item \textbf{Autorização por perfil}: confronta o perfil do utilizador (RF-01) com os requisitos da rota.
  \item \textbf{Validação de \emph{payload}}: valida o corpo do pedido contra o esquema Mongoose ou contra um \emph{schema} JSON.
  \item \textbf{Controlador da rota}: executa a lógica de negócio.
  \item \textbf{\emph{Middleware} de erros}: converte exceções em respostas HTTP coerentes com a OpenAPI.
\end{enumerate}


\subsection{Motor de regras e modo Automático}
\label{sec:motor-regras}

\begin{decisao}[title=Decisão DA-02: localização do motor de regras]
\textbf{Decisão.} O motor de regras corre na API, e não no cliente.

\textbf{Justificação.} Embora o cliente possa fornecer \emph{feedback}
otimista quando uma medição é registada \emph{offline}, a fonte da
verdade tem de estar no servidor: (i) os limites do plano podem
mudar; (ii) a coerência entre lotes exige uma visão global; (iii)
no modo \emph{Automático}, ações como ativar a rega ou ajustar a
ventilação têm efeitos físicos que não podem depender do estado
local do cliente.
\end{decisao}


\section{Modelo de dados}
\label{sec:modelo-dados}

% Apresentar o modelo conceptual com um diagrama (entidades e
% relações). Pode ser feito em TikZ ou desenhado num editor
% externo e incluído como imagem.

A Figura~\ref{fig:modelo-dados} apresenta o modelo conceptual
das entidades persistidas em MongoDB. Por simplicidade, são
omitidos atributos secundários.

% Substituir pelo diagrama real (e.g., draw.io exportado para PDF).
\begin{figure}[H]
\centering
\fbox{\parbox{0.9\linewidth}{\centering\vspace{2cm}
\textit{[Inserir aqui o diagrama de entidades e relações.]}\vspace{2cm}}}
\caption{Modelo conceptual de dados (entidades e relações).}
\label{fig:modelo-dados}
\end{figure}

\begin{nota}
Em MongoDB, cada documento pode ``embeber'' subdocumentos ou
referenciar outros documentos por \texttt{ObjectId}. As decisões
de \emph{embedding} versus \emph{referencing} influenciam o
desempenho das consultas e devem ser justificadas. Recomenda-se
uma subsecção dedicada onde, para cada par de entidades
relacionadas, se indique a opção tomada e o seu motivo.
\end{nota}


\section{Segurança}
\label{sec:seguranca}

% Reunir num só lugar todas as decisões de segurança.
% Útil para a equipa, útil para a avaliação.

As principais decisões de segurança adotadas são as seguintes:

\begin{itemize}
  \item As palavras-passe são armazenadas com \emph{hashing} adaptativo (\texttt{bcrypt}, fator de custo $\geq 10$).
  \item A autenticação usa JWT assinados (HS256 ou RS256) com \emph{payload} mínimo e expiração de [N] minutos.
  \item Todas as rotas, exceto as públicas, exigem JWT válido.
  \item A autorização é controlada por perfil; o servidor nunca confia em campos sensíveis enviados pelo cliente (e.g., o perfil do utilizador é sempre obtido do JWT, não do corpo do pedido).
  \item Os \emph{inputs} são validados contra esquemas; as saídas que possam conter texto introduzido pelo utilizador são \emph{escaped}.
  \item As ações sensíveis (criação de planos, autorização de planos pontuais, eliminação de utilizadores) ficam registadas em \texttt{LogAuditoria}.
\end{itemize}


\section{Resumo das decisões}
\label{sec:resumo-decisoes}

% Uma tabela síntese ajuda muito o avaliador a perceber o
% pensamento crítico da equipa. Listar as decisões principais
% com referência cruzada para a secção onde são justificadas.

\begin{table}[H]
\centering
\caption{Síntese das decisões de arquitetura.}
\label{tab:decisoes}
\renewcommand{\arraystretch}{1.3}
\begin{tabularx}{\linewidth}{l X l}
\toprule
\textbf{ID} & \textbf{Decisão} & \textbf{Secção} \\
\midrule
DA-01 & Armazenamento do JWT em [\dots]                              & \ref{sec:decisao-jwt} \\
DA-02 & Motor de regras na API                                       & \ref{sec:motor-regras} \\
DA-03 & Sincronização \emph{offline} via fila idempotente em IndexedDB & \ref{sec:sync} \\
DA-04 & [Adicionar conforme as decisões reais da equipa]             & --- \\
\bottomrule
\end{tabularx}
\end{table}

% =========================================================================
%  Capítulo 4 — Conclusão
% =========================================================================

\chapter{Conclusão}
\label{cap:conclusao}

% --- INSTRUÇÕES (apagar antes da entrega) ---------------------
% A conclusão NÃO é um resumo do que se fez. É uma reflexão
% crítica sobre:
%   - O que ficou bem feito.
%   - O que ficaria diferente se tivessem mais tempo.
%   - Que aprendizagens foram feitas.
%   - Que trabalho futuro faz sentido.
% Evitar frases vazias do tipo ``foi um projeto enriquecedor''.
% Concretizar: o que aprenderam SOBRE arquitetura web que não
% sabiam antes?
% ---------------------------------------------------------------


\section{Síntese do trabalho realizado}
\label{sec:sintese}

% Síntese curta do que foi efetivamente entregue. Esta secção
% deve ser objetiva: percentagem de requisitos implementados,
% módulos funcionais, etc.

[Apresentar de forma sintética o que foi efetivamente
implementado. Recomenda-se referir a cobertura dos requisitos
funcionais (RF-01 a RF-15) e dos requisitos não funcionais,
indicando claramente o que ficou completo, parcial e por
implementar.]


\section{Avaliação dos critérios de aceitação}
\label{sec:criterios}

% Confrontar o trabalho realizado com os critérios de aceitação
% do enunciado. Esta tabela facilita a vida ao avaliador.

A Tabela~\ref{tab:criterios} confronta o trabalho realizado com
os critérios de aceitação definidos no enunciado.

\begin{table}[H]
\centering
\caption{Verificação dos critérios de aceitação do projeto.}
\label{tab:criterios}
\footnotesize
\renewcommand{\arraystretch}{1.4}
\begin{tabularx}{\linewidth}{l X c}
\toprule
\textbf{ID} & \textbf{Critério} & \textbf{Estado} \\
\midrule
CA-01 & Um utilizador autenticado acede apenas às operações permitidas pelo seu perfil. & \\
CA-02 & O sistema suporta os três tipos de plano de cultivo e a sua associação a lotes. & \\
CA-03 & O registo de medições ambientais pode originar alertas automáticos. & \\
CA-04 & O sistema demonstra comportamento distinto em modo \emph{Manual} e \emph{Automático}. & \\
CA-05 & É possível consultar histórico, exportar relatórios e observar rastreabilidade. & \\
CA-06 & A arquitetura de armazenamento no \emph{browser} está implementada e coerente. & \\
CA-07 & A implementação está alinhada com a especificação OpenAPI. & \\
\bottomrule
\end{tabularx}
\end{table}

\noindent
\textit{Legenda:} preencher cada linha com \checkmark\ (cumprido),
$\sim$ (parcialmente cumprido) ou $\times$ (não cumprido), seguido
de uma breve evidência.


\section{Reflexão crítica}
\label{sec:reflexao}

% Esta é a secção que distingue um relatório bom de um relatório
% medíocre. Ser HONESTO. Os avaliadores valorizam mais um
% relatório que reconhece limitações do que um que pretende ser
% perfeito.

\subsection{O que correu bem}

[Identificar os aspetos do projeto onde a equipa considera ter
tomado boas decisões. Justificar com evidências concretas — não
basta dizer ``a arquitetura ficou bem''; é preciso explicar
porquê.]

\subsection{O que faríamos diferente}

[Identificar honestamente as decisões que, em retrospetiva, a
equipa repensaria. Esta é uma das partes mais valorizadas do
relatório: mostra capacidade de aprendizagem e pensamento
crítico.]

\subsection{Aprendizagens da equipa}

[Aprendizagens técnicas e de processo. Que conceitos ficaram mais
claros? Que ferramentas mostraram ser mais ou menos adequadas? Que
prática de equipa funcionou melhor?]


\section{Trabalho futuro}
\label{sec:futuro}

% Lista priorizada de melhorias possíveis. Pode incluir:
%   - Requisitos identificados mas não implementados.
%   - Otimizações de desempenho ou usabilidade.
%   - Funcionalidades que extravasam o âmbito atual mas
%     fariam sentido numa próxima iteração.

Identificam-se como direções relevantes para iterações futuras:

\begin{itemize}
  \item {}[Listar funcionalidades não implementadas que fariam sentido — e.g., notificações \emph{push}, \emph{dashboards} avançados, integração com sensores reais.]
  \item {}[Listar melhorias de natureza técnica — e.g., \emph{cobertura de testes automatizados, observabilidade, hardening de segurança}.]
  \item {}[Listar mudanças de processo ou metodologia que ajudariam num projeto semelhante.]
\end{itemize}


% --- Bibliografia ----------------------------------------------
\printbibliography[heading=bibintoc, title={Referências}]

% --- Anexos ----------------------------------------------------
\appendix
% =========================================================================
%  Anexos
% =========================================================================

\chapter{Especificação OpenAPI}
\label{anexo:openapi}

% Excerto da especificação OpenAPI ou referência ao ficheiro
% completo. Para excertos mais longos, considerar incluir como
% código com a package "listings".

A especificação completa da API está disponível no ficheiro
\texttt{openapi.yaml} no repositório do projeto
(\url{\repositorio}). A título ilustrativo apresenta-se a seguir
o excerto correspondente ao \emph{endpoint} de gestão de alertas
especificado na Sprint 1.

\begin{lstlisting}[language=, caption={Excerto de \texttt{openapi.yaml} — ignorar alerta com justificacao.}]
paths:
  /api/v1/alerts/{alertId}/ignore:
    post:
      summary: Ignorar alerta (justificacao obrigatoria)
      parameters:
        - in: path
          name: alertId
          required: true
          schema:
            type: string
      requestBody:
        required: true
        content:
          application/json:
            schema:
              type: object
              required: [reason]
              properties:
                reason: { type: string, minLength: 3 }
      responses:
        "200":
          description: Alerta ignorado
          content:
            application/json:
              schema:
                type: object
                properties:
                  id:
                    type: string
                  status:
                    type: string
                    enum: [ignorado]
                  ignoredReason:
                    type: string
        "400":
          description: Justificacao em falta ou invalida
\end{lstlisting}


\chapter{Dados de demonstração}
\label{anexo:dados}

% Descrição breve do dataset de demonstração entregue junto com
% o código (CSV ou Excel exemplo).

O dataset de demonstração inclui:

\begin{itemize}
  \item {}[N] utilizadores (um por perfil, no mínimo).
  \item {}[N] ervas aromáticas, importadas a partir do ficheiro \texttt{ervas-exemplo.csv}.
  \item {}[N] planos de cultivo (\emph{regular}, \emph{emergência} e \emph{pontual}).
  \item {}[N] lotes em diferentes estados (\emph{ativo}, \emph{concluído}, \emph{comprometido}).
  \item Histórico de medições ambientais e alertas associados.
\end{itemize}


\chapter{Distribuição do trabalho}
\label{anexo:distribuicao}

% Tabela honesta com a contribuição de cada elemento da equipa.
% Útil para a avaliação individual e para a equipa refletir.

\begin{table}[H]
\centering
\caption{Contribuição de cada elemento da equipa.}
\label{tab:distribuicao}
\renewcommand{\arraystretch}{1.4}
\begin{tabularx}{\linewidth}{l X c}
\toprule
\textbf{Elemento} & \textbf{Contribuições principais} & \textbf{\%} \\
\midrule
Aluno 1 & [Áreas de responsabilidade e principais entregas]   &      \\
Aluno 2 & [Áreas de responsabilidade e principais entregas]   &      \\
Aluno 3 & [Áreas de responsabilidade e principais entregas]   &      \\
\midrule
\textbf{Total} & & \textbf{100\%} \\
\bottomrule
\end{tabularx}
\end{table}


\end{document}
